%========     DECLARAÇÃO DO USO DE INTELIGÊNCIA ARTIFICIAL    =========%
\section*{\centering APÊNDICE A - DECLARAÇÃO DO USO DE INTELIGÊNCIA ARTIFICIAL}

\addcontentsline{toc}{section}{APÊNDICE A - DECLARAÇÃO DO USO DE INTELIGÊNCIA ARTIFICIAL}

\vspace{1cm}

Declaro, para os devidos fins, que o desenvolvimento deste trabalho contou com o suporte de ferramentas baseadas em Inteligência Artificial Generativa. O uso desses recursos ocorreu de forma assistida e complementar, com o propósito de acelerar e otimizar as etapas de pesquisa, desenvolvimento e redação.

O emprego das tecnologias distribuiu-se da seguinte forma:
\begin{itemize}
  \item \textbf{Claude}: Utilizado para a compreensão conceitual e estudos preliminares sobre o problema de transporte ótimo, além de auxiliar no processo de \textit{brainstorming} para a delimitação de escopo e restrições da modelagem;
  \item \textbf{Gemini}: Empregado como suporte na revisão bibliográfica e catalogação das fontes bibliográficas iniciais;
  \item \textbf{Assistente Antigravity (utilizando os modelos Gemini 3.5 Flash e Claude 4.6 Opus)}: Utilizado majoritariamente na escrita deste relatório, elaboração dos slides de apresentação, codificação da solução em Python e elucidação de tópicos específicos com base na bibliografia previamente selecionada.
\end{itemize}

Ressalta-se que todos os códigos, textos, equações e análises produzidos com o suporte dessas ferramentas foram revisados, validados e submetidos a um criterioso processo de avaliação e curadoria humana conduzido integralmente pelo autor. Desse modo, o recurso à Inteligência Artificial serviu estritamente como agente facilitador de produtividade científica, preservando a autoria, a responsabilidade acadêmica e o rigor conceitual do trabalho apresentado.

\vspace{1.5cm}

\noindent \textbf{Softwares e modelos utilizados}: Claude (Anthropic), Gemini (Google) e assistente Antigravity (com modelos Gemini 3.5 Flash e Claude 4.6 Opus).